%% =====================================================================
%%  B3-T-27-03 -- what B3 contributes to "The Siege Bill"
%%  -------------------------------------------------------------------
%%  LAYER A: APPEND-ONLY. Written once at the entry's write, then left
%%  alone. If the source has more to say later, add a bullet -- do not
%%  rewrite. Material found ONLY here goes in the `unique` slot with
%%  @unique set to yes, which makes it undeletable (rule R10).
%% =====================================================================
%% @kind: source
%% @id: B3-T-27-03
%% @book: B3
%% @topic: T-27-03
%% @locator: Law 18, pp. 129--138
%% @status: distilled
%% @unique: no
%% @integrated: yes
%% @updated: 2026-09-19

\dossier{B3}{T-27-03}{\bookshort{B3}}{Law 18, pp. 129--138}

\begin{distilled}
  \item The world is dangerous and enemies are everywhere, and the fortress seems safest --- but isolation exposes you to more dangers than it protects you from.
  \item Isolation cuts you off from valuable information, makes you conspicuous, and turns the citadel into a symbol of everything hateful in power that the citizens betray to the first enemy.
  \item Ch'in Shih Huang Ti walled himself in with mercury rivers and died believing himself immortal; the empire was rewritten within four years.
  \item Louis XIV's Versailles is the counter-model: place yourself at the centre of circulation, aware of everything happening on the street.
  \item Narrow reversal: isolation in small doses as a workshop for thought --- Machiavelli wrote The Prince in exile --- but the return road must stay open, and closed rooms grow strange ideas.
\end{distilled}
